\begin{table}[H]
\centering
\caption{Robustness of the adjusted college gap}
\label{tab:robustness}
\begin{threeparttable}
\small
\setlength{\tabcolsep}{4pt}
\begin{tabular}{lccccr}
\toprule
Specification & Pre-pandemic & Onset & Mid-pandemic & Post-pandemic & Obs. \\
\midrule
Baseline & -0.009$^{***}$ & -0.035$^{***}$ & 0.018$^{***}$ & -0.010$^{***}$ & 8,002,128 \\
 & (0.001) & (0.001) & (0.001) & (0.001) & \\
Prime age 25--54 & -0.008$^{***}$ & -0.037$^{***}$ & 0.009$^{***}$ & -0.012$^{***}$ & 5,570,238 \\
 & (0.001) & (0.001) & (0.001) & (0.001) & \\
Unweighted & -0.011$^{***}$ & -0.026$^{***}$ & 0.024$^{***}$ & -0.015$^{***}$ & 8,002,128 \\
 & (0.001) & (0.001) & (0.001) & (0.001) & \\
Base quarter 2019Q1 & -0.009$^{***}$ & -0.035$^{***}$ & 0.018$^{***}$ & -0.010$^{***}$ & 8,002,128 \\
 & (0.001) & (0.001) & (0.001) & (0.001) & \\
Outcome: informal employment in $t+1$ & -0.028$^{***}$ & -0.009$^{***}$ & -0.047$^{***}$ & -0.020$^{***}$ & 8,002,128 \\
 & (0.002) & (0.002) & (0.002) & (0.002) & \\
\bottomrule
\end{tabular}
\begin{tablenotes}[flushleft]\footnotesize
\item Notes: Each row reports survey-weighted average predictive differences in the outcome between workers with and without a college degree, from the full specification of equation~(\ref{eq:es}). The last row replaces the outcome by an indicator for being informally employed in $t+1$ -- a destination state, not a transition, so it includes workers who were already informal in $t$. Standard errors in parentheses are two-way clustered by primary sampling unit and year-quarter. Stars denote $^{*}p<0.10$, $^{**}p<0.05$, $^{***}p<0.01$. Periods are pre-pandemic 2012Q1--2019Q4, onset 2020Q1, mid-pandemic 2020Q2--2021Q3 and post-pandemic 2021Q4--2024Q4.
\end{tablenotes}
\end{threeparttable}
\end{table}
